\documentclass[11pt]{article}
\usepackage[hmargin=.88in,vmargin=.82in]{geometry}
\usepackage{amsmath,amssymb}
\usepackage[colorlinks=true,linkcolor=blue,citecolor=blue,urlcolor=blue,pdfauthor={Lennart Lopin},
 pdftitle={Smooth integers in multiplicative cosets at a squared-logarithmic cutoff}]{hyperref}
\setlength{\emergencystretch}{2em}
\newtheorem{theorem}{Theorem}
\newtheorem{lemma}{Lemma}
\newtheorem{corollary}{Corollary}
\newtheorem{proposition}{Proposition}
\newcommand{\cstar}{c_*}
\newcommand{\re}{\operatorname{Re}}
\title{Smooth integers in multiplicative cosets\\at a squared-logarithmic cutoff}
\author{Lennart Lopin}
\date{November 21, 2026}
\begin{document}
\maketitle
\begin{abstract}
Let \(\chi\) be a nonprincipal character modulo a prime \(p\) and assume
the Riemann hypothesis for \(L(s,\chi)\). Fix \(0<a<1\) and put
\(y=\eta(\log p)^2\), with \(\eta>\mathcal C(1)^{-2}\) in the
Fourier normalization of Carneiro, Milinovich, Quesada-Herrera and Ramos.
The character sum over \(y\)-smooth integers up to \(x\) is
\(O_{a,\eta}(\Psi(x,y)e^{-bu})\), uniformly in \(\chi\) and
\(p^a\le x\le p\), where \(u=\log x/\log y\) and \(b=b(a,\eta)>0\).
Under Dirichlet GRH, every coset of every subgroup of index
\(m\le e^{bu/2}\) therefore has its expected share with relative error
\(O_{a,\eta}(e^{-bu/2})\). In particular \(\eta=0.77\) is admissible.
The ordinary character sum satisfies the same bound for every fixed
\(\eta>2/e^2\) under RH for \(L(s,\chi)\).
\end{abstract}

\section{Character sums and coset counts}
Write \(P^+(n)\) for the greatest prime factor of \(n>1\), set \(P^+(1)=1\),
and let \(\Psi(x,y)\) count the \(n\le x\) with \(P^+(n)\le y\).
We write \(S(x,\chi)=\sum_{n\le x}\chi(n)\) and
\(S_y(x,\chi)=\sum_{n\le x,\,P^+(n)\le y}\chi(n)\).
All logarithms are natural. Write \(\mathrm{RH}(\chi)\) for the assertion
that every nontrivial zero of \(L(s,\chi)\) has real part \(1/2\). Put
\(\cstar=\mathcal C(1)^{-2}\) as in \cite[(1.1)]{cmqr}.

\begin{theorem}\label{th:main}
Fix \(0<a<1\) and \(\eta>\cstar\). There are positive constants \(b,K\)
and \(p_0\), depending only on \(a,\eta\), such that, for every prime
\(p\ge p_0\), every nonprincipal \(\chi\bmod p\) satisfying
\(\mathrm{RH}(\chi)\), and every real
\(p^a\le x\le p\),
\begin{equation}\label{eq:main}
 |S_y(x,\chi)|+|S(x,\chi)|\le K\Psi(x,y)e^{-bu},
 \qquad y=\eta(\log p)^2,\quad u=\frac{\log x}{\log y}.
\end{equation}
\end{theorem}
In particular the difference of the two sums satisfies the same bound.
The published bounds \(0.7596<\cstar<0.7615\) allow the cutoff \(0.77\)
\cite[Theorems 1 and 5]{cmqr}.

\begin{corollary}\label{cor:cosets}
Let \(B\le\mathbb F_p^\times\) have index \(m\) and suppose
\(\mathrm{RH}(\chi)\) holds for every nonprincipal character trivial
on \(B\). For every coset \(gB\), in the range of Theorem~\ref{th:main},
\[
 \#\{n\le x:P^+(n)\le y,\ n\bmod p\in gB\}
 =\frac{\Psi(x,y)}{m}+O_{a,\eta}(\Psi(x,y)e^{-bu}).
\]
The relative error is \(O_{a,\eta}(m e^{-bu})\), hence
\(O_{a,\eta}(e^{-bu/2})\) for \(m\le e^{bu/2}\). Under Dirichlet
GRH these estimates hold simultaneously for all such subgroups.
\end{corollary}
Indeed \(p>y\) for large \(p\), so all smooth \(n\) are units. The
indicator of \(gB\) on units equals
\(m^{-1}\sum_{\chi\in B^\perp}\overline{\chi(g)}\chi(n)\), where
\(B^\perp\) consists of the \(m\) characters trivial on \(B\).
The principal character contributes \(\Psi(x,y)/m\), and the other
\(m-1\) bounds are divided by the same \(m\).

For example, if \(p\equiv1\bmod3\), take a character \(\xi\) of order
three. The hypothesis \(\mathrm{RH}(\xi)\) suffices since
\(L(s,\bar\xi)=\overline{L(\bar s,\xi)}\). Write
\(\omega=e^{2\pi i/3}\). For \(j=0,1,2\), the number with
\(\xi(n)=\omega^j\) is exactly
\[
 \frac{\Psi(x,y)+2\re(\omega^{-j}S_y(x,\xi))}{3}
 =\frac{\Psi(x,y)}3+O_{a,\eta}(\Psi(x,y)e^{-bu}).
\]
This includes the cubic residues (\(j=0\)) and distinguishes the two
other cosets, including at \(x=\sqrt p\) with \(a=1/2\).

\paragraph{Relation to prior estimates.}
Granville and Soundararajan conjectured a uniform approximation to
ordinary character sums by their smooth restrictions. Their Theorem~2
proves \(S(x,\chi)-S_y(x,\chi)\ll\Psi(x,y)/(\log\log q)^2\) under
RH for \(L(s,\chi)\) when \(1\le x\le q\) and
\(y\ge(\log q)^2(\log x)^2(\log\log q)^{12}\)
\cite[Conjecture 1 and Theorem 2]{gs}.

Dong--Zhao--Zhou \cite[Theorem 1.4, v1]{dzz} assume RH for
\(L(s,\chi)\) and cover all nonprincipal characters of arbitrary
modulus \(q\), including imprimitive characters. They bound both
sums relative to the principal smooth count over \(q^a\le x\le q\).
Their cutoff is \((\log q+(\log x)^2)(\log\log q)^C\) with \(C>6\).
Their error factor is
\(\exp(-\gamma(\log x)\log_3 q/\log_2 q)\) for \(0<\gamma<(C-6)/4\),
where \(\log_k\) denotes the \(k\)-fold iterated logarithm. At
\(q=p,\ x=p^\theta\) the ratio of their cutoff to ours is
\((\theta^2+1/\log p)(\log\log p)^C/\eta\), tending uniformly to
infinity on \([a,1]\). Their decay exponent includes a further factor
of order \(\log_3 p\). At prime modulus their bound also gives
equidistribution for growing subgroup indices by the orthogonality
argument above, at their cutoff and under the corresponding character
hypotheses. Our estimate gives this conclusion at a fixed multiple
of \((\log p)^2\). Decreasing the cutoff removes a signed subsum
whose size is not controlled by a bound for the original sum.

Croot--Granville--Pemantle--Tetali already study smooth squares and
compare the two sums under GRH \cite[\S7.3, Lemma 7]{cgpt}. At our
cutoff their error factor \((\log p)^3/\sqrt y\) equals
\((\log p)^2/\sqrt\eta\). Soundararajan uses a saddle-point weighted
distance to \(\ell^{it}\) for smooth Euler products in a range of moduli
bounded by a fixed power of \(y\) \cite[\S4]{sound}.
Lagarias and Soundararajan use the explicit formula for
\(\sum_{n\le v}\Lambda(n)\chi(n)n^{-it}\) to bound the truncated Euler
product \(L(s,\chi;y)\). Their Proposition~5.1 gives
\(|L(s,\chi;y)|\ll_\varepsilon(q|s|)^\varepsilon\) under GRH for
nonprincipal primitive \(\chi\bmod q\) and
\(\re s\ge1/2+\varepsilon\), with \(\varepsilon>0\) fixed
\cite[\S5]{ls}. Here the saddle satisfies
\(\alpha-1/2\asymp_{a,\eta}1/\log y\) by \eqref{eq:saddle}.
We need a saving relative to the positive Euler product on that line,
which the translated Fourier estimate gives in \eqref{eq:Us}.
Harper's individual
bounds impose conductor \(r\le y^{b_0}\) or \(r\le x^{d_0}\) for small
absolute constants \(b_0,d_0>0\) \cite[Theorem 3 and Proposition 1]{harper}.
Here \(r=p\ge x\).

\section{The translated explicit formula}
Put \(L=\log p,\ H=\log y,\ X=\log x\), so \(u=X/H\). Throughout,
\(\ell\) denotes a prime. Use the Fourier transform with kernel
\(e^{-2\pi i z v}\).

\begin{lemma}\label{lem:distance}
For every fixed \(\eta>\cstar\) there exist \(0<d<1\) and \(\delta>0\),
depending only on \(\eta\), such that for every sufficiently large
prime \(p\), every nonprincipal \(\chi\bmod p\) satisfying
\(\mathrm{RH}(\chi)\), and every
\(|t|\le\exp(L/H)\),
\begin{equation}\label{eq:distance}
 \sum_{dy<\ell\le y}\bigl(1-\re(\chi(\ell)\ell^{-it})\bigr)
 \ge \delta\,y/H .
\end{equation}
\end{lemma}
\emph{Proof.}
Choose a test \(F\) with real \(\widehat F\in C_c^\infty([-N,N])\) and
\[
 D(F)=\int_{-\infty}^0\widehat F(v)e^{\pi v}\,dv
      -\int_0^\infty |\widehat F(v)|e^{\pi v}\,dv,\qquad
 \mu_\eta=\sqrt\eta\,D(F)-\|F\|_1/(2\pi)>0.
\]
Such a choice follows from \cite[Theorem 4(i), (4.13) and \S4.4]{cmqr}.
Set \(d=e^{-2\pi N}\), \(\Delta=H/(2\pi)>N\), and form their test
\[
 \widehat h_0(v)=\tfrac12\{\widehat F(v-\Delta)+\widehat F(-v-\Delta)\},
 \qquad h_t(z)=h_0(z-t).
\]
On the real axis \(h_0(z)=\re(e^{2\pi i\Delta z}F(z))\), since
\(\widehat F\) is real. Thus \(|h_t(z)|\le |F(z-t)|\) and
\(\widehat h_t(v)=e^{-2\pi itv}\widehat h_0(v)\).
The analytic continuation required in \cite[(3.1)]{cmqr} exists for each
test by Fourier inversion on compact support. We now track the dependence
on the translation on the real axis.

For ordinates of \(L(s,\chi)\), counted with multiplicity, we claim
\begin{equation}\label{eq:zeros}
 \sum_{\gamma_\chi}|F(\gamma_\chi-t)|
 =\frac{L}{2\pi}\|F\|_1+
 O_F\!\left(\frac{L}{\log L}+\log(3+|t|)\right).
\end{equation}
To verify this for complex characters, take the signed cumulative count
\(N(v,\chi)\) for \(v>0\) and \(-N(-v,\bar\chi)\) for \(v<0\),
with the endpoint convention of \cite[(3.3)]{cmqr}. The negative
ordinates of \(\chi\) reflect to the positive ordinates of \(\bar\chi\), and
the two halves at zero give the entire central multiplicity. The
distributional derivative of this count is the zero-counting measure.
By \cite[(3.3)--(3.4)]{cmqr}, subtracting
\(M(v)=v\log(p|v|/(2\pi e))/(2\pi)\), \(M(0)=0\), leaves a remainder
\[
 R_\chi(v)\ll L/\log L+\log(3+|v|)
\]
on both half-lines, with an absolute implied constant. The bound on the
argument in (3.4) holds under RH for the individual primitive character
\cite[Theorem 5 and Example 8]{ccm}. Conjugation supplies the same
hypothesis for \(\bar\chi\). Half values at
jumps have no effect on integration by parts. Integrate
\(|F(v-t)|\) against \(dM+dR_\chi\).
In \(dM\) the \(L/(2\pi)\) density gives the main term.
The part \(\log(|v|/(2\pi))\) costs \(O_F(\log(3+|t|))\), using the
integrability of \(|\log|v||\) near \(0\) and the decay of \(F\).
For \(dR_\chi\), use absolute continuity of \(|F|\),
\(|(|F|)'|\le |F'|\) almost everywhere, and
\(\log(3+|v|)\le\log(3+|t|)+\log(3+|v-t|)\).
This proves \eqref{eq:zeros} without differentiating the modulation
\(e^{2\pi i\Delta z}\).

Every nonprincipal character modulo \(p\) is primitive. In (3.1) the
conductor term vanishes since \(\widehat h_t(0)=0\). The gamma integral
is \(O_F(\log(3+|t|))\) by Stirling's formula for either parity and
\(|h_t(v)|\le |F(v-t)|\). At \(v=\log n/(2\pi)>0\), the two prime
terms contain \(\chi(n)n^{-it}\) and its conjugate. Using
\(\widehat h_0(v)=\widehat F(v-\Delta)/2\) there gives
\begin{align}
 W_\chi(t)&=\frac1{2\pi}\sum_{n\ge2}\frac{\Lambda(n)}{\sqrt n}
    \widehat F\!\left(\frac{\log(n/y)}{2\pi}\right)
    \re(\chi(n)n^{-it}),\nonumber\\
 |W_\chi(t)|&\le\frac{L}{2\pi}\|F\|_1+
 O_F\!\left(L/\log L+\log(3+|t|)\right).\label{eq:W}
\end{align}
In particular the error is \(o_\eta(L)\) up to \(\exp(L/H)\), uniformly
over character order. The factor \(1/(2\pi)\) includes both prime
terms and the half in the definition of \(\widehat h_0\).

Write \(w_\ell=(\log\ell)\widehat F(\log(\ell/y)/(2\pi))/
(2\pi\sqrt\ell)\). Prime powers of degree at least two in
\eqref{eq:W} cost \(O_F(H^2)\), by counting them in \([dy,y/d]\).
Since \(w_\ell=0\) at the support endpoints, a lower bound for \(W_\chi(t)\)
is
\[
 \sum_{dy<\ell\le y} w_\ell-\sum_{y<\ell\le y/d}|w_\ell|
 - C_F\frac{H}{\sqrt y}
       \sum_{dy<\ell\le y}(1-\re(\chi(\ell)\ell^{-it}))-O_F(H^2).
\]
Indeed \(w_\ell r\ge w_\ell-|w_\ell|(1-r)\) for \(-1\le r\le1\).
The first two sums together equal \(\sqrt y\,D(F)+o_F(\sqrt y)\) by
the unconditional prime number theorem and the substitution \(\ell=ye^{2\pi v}\).
Compared with \eqref{eq:W} this leaves \(\mu_\eta L+o_\eta(L)\) for
the loss term. Thus \eqref{eq:distance} holds, for example with
\(\delta=\mu_\eta/(2C_F\sqrt\eta)\) for large enough \(p\).
\hfill\(\square\)

\section{The saddle point and the sharp smooth sum}
For \(f=1,\chi\) let \(Z_f(s)=\prod_{\ell\le y}(1-f(\ell)\ell^{-s})^{-1}\).
Let \(\alpha>0\) be the unique solution of
\(\sum_{\ell\le y}\log\ell/(\ell^\alpha-1)=X\), and put
\(Z=Z_1(\alpha)\), \(\theta=X/L\). Uniformly for \(a\le\theta\le1\),
\begin{align}
 \alpha&=\tfrac12+
      \{\log(2\sqrt\eta/\theta)+o_{a,\eta}(1)\}/H,
 & y^{-\alpha}&\sim X/(2y),\label{eq:saddle}\\
 \frac{x^\alpha Z}{\Psi(x,y)}&\asymp_{a,\eta}\sqrt{XH},
 &\log\Psi(x,y)&=\tfrac12 X+
   \{1+\log(2\sqrt\eta/\theta)+o_{a,\eta}(1)\}u.\label{eq:psi}
\end{align}
Indeed at \(\sigma=1/2+v/H\) the saddle sum is
\((2e^{-v}+o(1))\sqrt y\) by the prime number theorem, uniformly on
compact \(v\)-intervals. Higher powers in the sum cost \(O(H)\) since
\(\ell^{-2\sigma}\ll1/\ell\). Bracket \([a,1]\) by two fixed values
of \(2\sqrt\eta e^{-v}\), then invert by monotonicity.
Partial summation at this saddle gives
\(\log Z\sim y^{1-\alpha}/((1-\alpha)H)\sim u\) and
\((\log Z_1)''(\alpha)\sim XH\). In these two estimates the higher
powers cost \(O(\log H)\) and \(O(H^2)\), respectively.
The Hildebrand--Tenenbaum saddle-point formula gives \eqref{eq:psi}
\cite[Theorem 1.6]{ht}.

Put \(U_f(t)=|Z_f(\alpha+it)|/Z\). For the complex character the
logarithm of the product gives
\[
 -\log U_\chi(t)=
 \sum_{\ell\le y}\sum_{k\ge1}
 \frac{1-\re((\chi(\ell)\ell^{-it})^k)}{k\ell^{k\alpha}}.
\]
All summands are nonnegative. Restricting to \(k=1\) and \((dy,y]\)
in Lemma~\ref{lem:distance} gives, for some \(k_1>0\),
\begin{equation}\label{eq:Us}
 U_\chi(t)\le e^{-k_1u}\qquad (|t|\le e^{L/H}),
\end{equation}
since \((y/H)y^{-\alpha}\sim u/2\).

For \(0<\kappa<1\) fixed put \(T=e^{\kappa u}\), so
\(2\le T\le e^{L/H}\) for all sufficiently large \(p\).
Gorodetsky's proof of Lemma 3.5 \cite[(3.10)--(3.16)]{gor} gives
\begin{equation}\label{eq:perron}
 S_y(x,\chi)=\frac1{2\pi i}\int_{\alpha-iT}^{\alpha+iT}
 Z_\chi(s)\frac{x^s}{s}\,ds+
 O\!\left(\frac{x^\alpha Z}{T}
       +\frac{\Psi(x,y)\log T}{T^\alpha}+1\right).
\end{equation}
Indeed the error in truncated Perron is majorized, up to an absolute
constant, by
\[
 x^\alpha\sum_{P^+(n)\le y}
 \frac{|\chi(n)|}{n^\alpha(1+T|\log(x/n)|)}.
\]
Dropping \(|\chi(n)|\le1\) gives exactly the positive sum bounded in
the proof of that lemma. The \(O(1)\) also allows the full weight at
an integer endpoint. The lemma is unconditional and requires
\(x\ge y\ge C,\ T\ge2,\ 0<\alpha\le1\), which hold uniformly here by
\eqref{eq:saddle}. By \eqref{eq:Us} the integral in \eqref{eq:perron}
is \(O(x^\alpha Z e^{-k_1u}\log(2T))\). On dividing by \(\Psi(x,y)\),
\eqref{eq:psi} therefore bounds the absolute ratio by
\[
 \ll_{a,\eta}\sqrt{XH}\bigl(e^{-k_1u}\log(2T)+T^{-1}\bigr)
       +T^{-\alpha}\log T+\Psi(x,y)^{-1}.
\]
Here \(\alpha\to1/2\), \(\log\Psi(x,y)\sim X/2\), and
\(\log L=o(u)\) uniformly since \(u\ge aL/H\). Each term is thus
\(O_{a,\eta}(e^{-b_1u})\) for some \(b_1>0\), also at \(x=p^a,p\).
This proves the smooth-sum estimate.

\section{The ordinary character sum}
\begin{proposition}\label{prop:ordinary}
Fix \(0<a<1\) and \(\eta>2/e^2\). For every sufficiently large prime
\(p\), every nonprincipal \(\chi\bmod p\) satisfying \(\mathrm{RH}(\chi)\),
and \(p^a\le x\le p\), one has
\[
 |S(x,\chi)|\ll_{a,\eta}\Psi(x,y)e^{-b_2u},
 \qquad y=\eta(\log p)^2,\quad u=\log x/\log y ,
\]
with \(b_2=b_2(a,\eta)>0\) and the threshold for \(p\) depending only
on \(a,\eta\).
\end{proposition}
\emph{Proof.}
Let \(J=\log L\), \(V=\sqrt{x}\log x\), and
\(\sigma=1/2+\lambda/(2J)\) with \(\lambda\) in a fixed compact
subinterval of \((0,\infty)\). For \(|t|\le V\) the degree-one analytic
conductor \(Q(t)=p(3+|it+\mu|)\), \(\mu\in\{0,1\}\), satisfies
\[
 L\le\log Q(t)\le (1+\theta/2)L+O(J),\qquad
 \log\log Q(t)=J+O(1).
\]
Chirre's Theorem 1(i) under \(\mathrm{RH}(\chi)\) gives
\begin{equation}\label{eq:chirre}
 \log|L(\sigma+it,\chi)|
 \le\frac{(1+\theta/2)L}{2J}\log(1+e^{-\lambda})+O(L/J^2).
\end{equation}
Indeed his first term is
\(\log Q(t)\log(1+(\log Q(t))^{1-2\sigma})/(2\log\log Q(t))\).
His term involving \((2\sigma-1)/(\sigma(1-\sigma))\) and his upper
error are \(O(L/J^2)\) uniformly on the chosen compact interval.
The hypothesis \((1-\sigma)^2\log\log Q(t)\ge1\) holds for large \(p\)
\cite[Theorem 1(i)]{chirre}.

Average \(S(xe^{hW},\chi)\) over a standard normal \(W\), with
\(h=x^{-1/2}\). Its distance from \(S(x,\chi)\) is
\(O(1+x\,\mathbb E|e^{hW}-1|)=O(\sqrt x)\).
Averaging Perron's kernel on \(\re s=2\) inserts \(e^{s^2/(2x)}\).
Shift the integral to \(\re s=\sigma\). Partial summation from the
periodic bound \(|S(z,\chi)|\le p\) gives
\(|L(s,\chi)/s|\le p/\re s\) throughout this strip, justifying the
shift. The two tails \(|t|>V\) on the new line cost
\[
 \ll p x^\sigma \frac{x}{V}e^{-V^2/(2x)}
 \ll_{a,A}x^{-A}\qquad(A>0)
\]
since \(V^2/x=(\log x)^2\) and \(p\le x^{1/a}\). On \([-V,V]\) the
absolute kernel after removing \(x^\sigma\) has integral \(O(\log(2V))\).
Using \eqref{eq:chirre} and \(H=2J+\log\eta\), we obtain
\[
 |S(x,\chi)|\ll\sqrt x+
 \sqrt x\exp\!\left\{\left(
 \lambda+(1/\theta+1/2)\log(1+e^{-\lambda})+o(1)\right)u\right\}.
\]

Take \(\lambda=\log(1/\theta)+s_0\) for a fixed \(s_0>0\).
Then the exponent coefficient is at most
\(\log(1/\theta)+s_0+g(\theta)+o(1)\), where
\(g(\theta)=(1/\theta+1/2)\log(1+\theta)\).
This function attains its maximum on \([a,1]\) at \(1\). In fact
\[
 \theta^2 g'(\theta)=\frac{w-w^{-1}}2-\log w,\qquad w=1+\theta .
\]
The right side vanishes at \(w=1\) and has derivative
\((w-1)^2/(2w^2)\ge0\). Hence \(g(\theta)\le(3/2)\log2\).
By \eqref{eq:psi} the margin in the exponent is at least
\[
 1+\log(2\sqrt\eta)-\tfrac32\log2-s_0+o_{a,\eta,s_0}(1).
\]
For \(\eta>2/e^2\) choose \(0<s_0<
(1+\log(2\sqrt\eta)-\tfrac32\log2)/2\).
Now \(\lambda\in[s_0,\log(1/a)+s_0]\) and the saving is uniform.
The \(O(\sqrt x)\) error is absorbed by the same margin. \hfill\(\square\)

Since \(2/e^2<0.7596<\cstar\), Proposition~\ref{prop:ordinary} and
the smooth-sum estimate give Theorem~\ref{th:main} with
\(b=\min(b_1,b_2)\). The smooth sum still requires \(\eta>\cstar\)
in Lemma~\ref{lem:distance}, leaving the growth-rate question for the
least quadratic nonresidue open.

\small
\begin{thebibliography}{99}
\raggedright
\setlength{\itemsep}{2pt}
\setlength{\parskip}{0pt}
\bibitem{ccm} E.~Carneiro, V.~Chandee and M.~B.~Milinovich,
\emph{A note on the zeros of zeta and \(L\)-functions}.
\href{https://arxiv.org/pdf/1503.00955}{arXiv:1503.00955}.
\bibitem{cmqr} E.~Carneiro, M.~B.~Milinovich, E.~Quesada-Herrera and
A.~P.~Ramos, \emph{Fourier optimization, the least quadratic non-residue,
and the least prime in an arithmetic progression}.
\href{https://arxiv.org/pdf/2404.08380v2}{arXiv:2404.08380v2}.
\bibitem{chirre} A.~Chirre, \emph{A note on entire \(L\)-functions}.
\href{https://arxiv.org/pdf/1805.01119}{arXiv:1805.01119}.
\bibitem{cgpt} E.~Croot, A.~Granville, R.~Pemantle and P.~Tetali,
\emph{Running time predictions for factoring algorithms}, ANTS VIII
(2008). \href{https://www2.math.upenn.edu/~pemantle/papers/square-algorithms.pdf}{Authors' version}.
\bibitem{dzz} Z.~Dong, S.~Zhao and G.-L.~Zhou,
\emph{Smooth approximations to character sums and zeta sums}.
\href{https://arxiv.org/pdf/2611.11389v1}{arXiv:2611.11389v1}.
\bibitem{gor} O.~Gorodetsky, \emph{Smooth integers and de Bruijn's
approximation \(\Lambda\)}, Proc.\ Roy.\ Soc.\ Edinburgh Sect.~A
\textbf{155} (2025), 792--820.
\href{https://doi.org/10.1017/prm.2023.115}{doi:10.1017/prm.2023.115}.
\bibitem{gs} A.~Granville and K.~Soundararajan, \emph{Large character
sums}, J.~Amer.\ Math.\ Soc.\ \textbf{14} (2001), 365--397.
\href{https://dms.umontreal.ca/~andrew/PDF/ChiSums.pdf}{Authors' version}.
\bibitem{harper} A.~J.~Harper, \emph{Bombieri--Vinogradov and
Barban--Davenport--Halberstam type theorems for smooth numbers}.
\href{https://arxiv.org/pdf/1208.5992}{arXiv:1208.5992}.
\bibitem{ht} A.~Hildebrand and G.~Tenenbaum, \emph{Integers without
large prime factors}, J.~Th\'eor.\ Nombres Bordeaux \textbf{5} (1993),
411--484. \href{https://doi.org/10.5802/jtnb.101}{doi:10.5802/jtnb.101}.
\bibitem{ls} J.~C.~Lagarias and K.~Soundararajan, \emph{Counting smooth
solutions to the equation \(A+B=C\)}, Proc.\ London Math.\ Soc.\ (3)
\textbf{104} (2012), 770--798.
\href{https://arxiv.org/pdf/1102.4911v2}{arXiv:1102.4911v2}.
\bibitem{sound} K.~Soundararajan, \emph{The distribution of smooth
numbers in arithmetic progressions}.
\href{https://arxiv.org/pdf/0707.0299}{arXiv:0707.0299}.
\end{thebibliography}
\end{document}
